\section*{Chapter \ref{ch:g7:markets-and-prices} --- Markets and Prices}

\begin{solution}{exo:g7:markets-and-prices:1}
About 62 million tonnes in 2009 and about 42 million tonnes in 2010 (the table gives $61.7$
and $41.5$).
\end{solution}

\begin{solution}{exo:g7:markets-and-prices:2}
A \omterm{def:g7:markets-and-prices:shortage}{shortage} of $120 - 80 = 40$ kg.
\end{solution}

\begin{solution}{exo:g7:markets-and-prices:3}
A \omterm{def:g7:markets-and-prices:market-surplus}{market surplus} of $90 - 50 = 40$ kg.
\end{solution}

\begin{solution}{exo:g7:markets-and-prices:4}
It fell, from $4.82$ dollars to $2.04$ dollars a dozen.
\end{solution}

\begin{solution}{exo:g7:markets-and-prices:5}
At the old \omterm{def:g7:money:price}{price}, buyers would want more strawberries than are now offered: there is a
\omterm{def:g7:markets-and-prices:shortage}{shortage}, and the \omterm{def:g7:money:price}{price} rises. The higher \omterm{def:g7:money:price}{price} leads buyers to buy fewer strawberries and
sellers to bring some from other regions. The \omterm{def:g7:money:price}{price} is higher and the quantity sold smaller
than before the frost.
\end{solution}

\begin{solution}{exo:g7:markets-and-prices:6}
At the old \omterm{def:g7:money:price}{price}, more apples are offered than buyers want: there is a \omterm{def:g7:markets-and-prices:market-surplus}{market surplus}, and
sellers lower their \omterm{def:g7:money:price}{prices}. The lower \omterm{def:g7:money:price}{price} leads buyers to buy more. The \omterm{def:g7:money:price}{price} is lower and
the quantity sold larger.
\end{solution}

\begin{solution}{exo:g7:markets-and-prices:7}
$(5.37 - 0.89) \div 5.37 = 4.48 \div 5.37 \approx 0.834$: a fall of $83.4\,\%$.
\end{solution}

\begin{solution}{exo:g7:markets-and-prices:8}
(a) $684.7 - 653.7 = 31.0$ million tonnes less. (b) $698.0 - 653.7 = 44.3$ million tonnes
more. (c) The part where a higher \omterm{def:g7:money:price}{price} tells growers to bring more: the world's wheat area
grew in 2011. (The weather of 2011 also mattered; the numbers alone cannot separate the two.)
\end{solution}

\begin{solution}{exo:g7:markets-and-prices:9}
For example: a \omterm{def:g7:markets-and-prices:market}{market} where the \omterm{def:g7:money:price}{price} is free to move, a harvest that fails, and a \omterm{def:g7:money:price}{price} that
falls although nothing else has changed (no new supply, no fall in what buyers want).
\end{solution}

\begin{solution}{exo:g7:markets-and-prices:10}
(a) $9\,000 - 6\,000 = 3\,000$ kg a week. (b) $3\,000 \div 3 = 1\,000$ households. (c) For
example: queue earlier and longer, send several members of the family, or buy potatoes to
resell them above 1 dalen to those who did not get any.
\end{solution}

\begin{solution}{exo:g7:markets-and-prices:11}
(a) Four buyers are willing to pay 25 dalens or more, and the fifth only 24: the \omterm{def:g7:money:price}{price} ends
above 24 and at most 25 dalens. (b) With 6 boxes, the sixth buyer pays up to 22 and the
seventh only 20: the \omterm{def:g7:money:price}{price} ends above 20 and at most 22 dalens.
\end{solution}

\begin{solution}{exo:g7:markets-and-prices:12}
For example: did poorer households pay more of their budget than richer ones? How much did
households cut the eggs they bought? Did \omterm{def:g7:money:price}{prices} rise more in some regions? A survey of
household spending, with the eggs bought by each household and its income, would answer the
first.
\end{solution}

\begin{solution}{pb:g7:markets-and-prices:1}
\textbf{1.} A \omterm{def:g7:markets-and-prices:shortage}{shortage}: $40 - 20 = 20$ tonnes.
\textbf{2.} A \omterm{def:g7:markets-and-prices:market-surplus}{market surplus}: $50 - 10 = 40$ tonnes.
\textbf{3.} 3 dalens per kilogram; 30 tonnes a week.
\textbf{4.} $30 \times 1\,000 = 30\,000$ kg at 3 dalens: $90\,000$ dalens a week.
\textbf{5.} 5, 10, 15, 20 and 25 tonnes at 1, 2, 3, 4 and 5 dalens.
\textbf{6.} At 4 dalens, 20 tonnes are wanted and 20 offered: the \omterm{def:g7:money:price}{price} is 4 dalens and 20
tonnes are sold.
\textbf{7.} \omterm{def:g7:money:price}{Price}: $(4 - 3) \div 3 \approx 0.333$, a rise of $33.3\,\%$. Quantity: $(30 - 20)
\div 30 \approx 0.333$, a fall of $33.3\,\%$.
\textbf{8.} The higher \omterm{def:g7:money:price}{price} tells households to buy fewer apples (or other fruit instead), and
tells orchards elsewhere and stores of apples that it pays to sell in Brenmoor.
\textbf{9.} 30 tonnes wanted, 15 tonnes offered.
\textbf{10.} $30 - 15 = 15$ tonnes, that is $15\,000$ kg.
\textbf{11.} $15\,000 \div 2 = 7\,500$ households.
\textbf{12.} Those who arrive first, or who know the sellers. Queues form, and some people resell
apples above 3 dalens; orchards may send apples to towns where the \omterm{def:g7:money:price}{price} is free.
\end{solution}
